\section{Proposed Adaptive Extension}
\label{sec:method}

\textit{This section specifies a proposed mechanism; it was not implemented or
evaluated in the current study.}  The fixed executor in Sec.~\ref{sec:execution}
tests whether different groups can retain predictions for different durations.
The next step is to make those durations depend on the current prediction rather
than fixing them per task.

\subsection{Group-wise temporal reliability}

At a policy query, let the frozen policy predict a $K$-step chunk
\begin{equation}
\widehat{\mathbf A}_t=\pi_\theta(o_t)
 = (\widehat{\mathbf a}_{t\mid t},\ldots,
    \widehat{\mathbf a}_{t+K-1\mid t}).
\label{eq:method_chunk}
\end{equation}
For action group $g$ and future offset $k$, we define a temporal-reliability
quantity
\begin{equation}
R_g(k\mid o_t,\widehat{\mathbf A}_t)
 = \Pr\!\left[\mathcal V_g(t,k)=1\mid o_t,\widehat{\mathbf A}_t\right],
\label{eq:reliability}
\end{equation}
where $\mathcal V_g(t,k)$ denotes that the group-$g$ prediction at offset $k$
remains valid under a predeclared group-wise criterion.  The criterion is not
assumed to be a task-success label: it can measure agreement with a future
target, consistency with a fresh policy query, or calibrated predictive
uncertainty.

The proposed horizon rule is
\begin{equation}
h_g(t)=\max\left\{k\in\{1,\ldots,K\}:
R_g(k\mid o_t,\widehat{\mathbf A}_t)>\tau\right\},
\label{eq:adaptive_horizon}
\end{equation}
with a reliability threshold $\tau$.  In a deployment implementation, the
resulting $h_g(t)$ would control how long group $g$ retains its accepted
subsequence.  It would not create an independent policy-query stream: if any
group expires, the frozen policy is queried once for a new full joint chunk, and
only expired groups accept their slices, exactly as in Sec.~\ref{sec:execution}.
Thus the proposed rule changes group acceptance horizons, not the policy's
observation or query semantics.

\subsection{Candidate self-supervised targets}

No target or estimator architecture is selected in this draft.  The following
choices are concrete hypotheses for the next study:
\begin{enumerate}
    \item \textbf{Future action consistency.}  On demonstration data, label
    $\mathcal V_g(t,k)$ when the predicted group action agrees with the
    demonstrated action at the same offset under a predeclared distance and
    tolerance.  This is directly self-supervised but does not include
    closed-loop success labels.
    \item \textbf{Policy re-query disagreement.}  Compare the stored prediction
    with the corresponding slice of a fresh prediction at a later observation.
    This targets stale predictions at deployment, but requires an additional
    policy query during training or calibration.
    \item \textbf{Predictive uncertainty.}  Estimate a group-wise distribution
    or dispersion from an ensemble, sampled action chunks, or model-specific
    denoising statistics, and calibrate it to the validity event.
\end{enumerate}

These choices may produce nonmonotonic, task-dependent reliability curves, as
suggested by the offline persistence audits in Sec.~\ref{sec:results}.  The
current paper therefore treats Eq.~\ref{eq:adaptive_horizon} as a proposed
adaptive execution mechanism and reports no learned-estimator accuracy,
closed-loop dynamic horizon result, or causal benefit from it.
